\documentclass[../../../main.tex]{subfiles}
\begin{document}

\subsection{Pooling Effect Sizes}

\subsubsection{Fixed-Effect Model}

The fixed-effect model assumes that all effect sizes stem from a single, homogeneous population. It states that all studies share the same true effect size. According to the fixed-effect model, the only reason why a study $k$'s observed effect size $\hat{\theta}_k$ deviates from $\theta$ is because of its sampling error $\epsilon_k$. This sampling error causes the observed effect to deviate from the overall, true effect. We can describe the relationship as the following \autocite{Harrer2021}.

\begin{equation}
\hat{\theta}_k = \theta + \epsilon_k.
\end{equation}

In the formula of the fixed-effect model, the true effect size is symbolized by $\theta$, not $\theta_k$ as in \autoref{Equation:SampleError}. Previously, we only made statements about the true effect size of one individual study $k$. The fixed-effect tells us that a study's true effect size $\theta_k$, and the overall, pooled effect size $\theta$, are identical. The formula of the fixed-effect models tells us that there is only one reason why observed effect sizes $\theta_k$ deviate from the true overall effect: because of the sampling error $\epsilon_k$. Furthermore, all things being equal, as the sample size becomes larger, the sampling error becomes smaller \autocite{Harrer2021}.

If we want to calculate the pooled effect size under the fixed-effect model, we therefore simply use a weighted average of all studies. To calculate the weight $w_k$ for each study $k$, we can use the standard error, which we square to obtain the variance $s_k^2$ of each effect size. Since a lower variance indicates higher precision, the inverse of the variance is used to determine the weight of each study \autocite{Harrer2021}.

\begin{equation}
w_k = \frac{1}{s_k^2}.
\end{equation}

Once we know the weights, we can calculate the weighted average, our estimate of the true pooled effect $\hat{\theta}$. This method is the most common approach to calculate average effects in meta-analyses. Since we use the inverse of the variance, it is often called inverse-variance weighting, or simply inverse-variance meta-analysis \autocite{Harrer2021}.

\begin{equation}
\hat{\theta} = \frac{\sum_{k=1}^{K}\hat{\theta}_kw_k}{\sum_{k=1}^{K}w_k}.
\end{equation}

\subsubsection{Random-Effects Model}

The fixed-effect model assumes that all our studies are part of a homogeneous population. However, it is simply unrealistic that studies in a meta-analysis are always completely homogeneous. It is likely that we can anticipate considerable between-study heterogeneity in the true effects. The random-effects model assumes that there is not only one true effect size but a distribution of true effect sizes. The goal of the random-effects model is therefore not to estimate the one true effect size of all studies, but the mean of the distribution of true effects \autocite{Harrer2021}.

Similar to the fixed-effect model, the random-effects model starts by assuming that an observed effect size $\hat{\theta}_k$ is an estimator of the study's true effect size $\theta_k$, burdened by sampling error $\epsilon_k$ \autocite{Harrer2021}.

\begin{equation}
\hat{\theta}_k = \theta_k + \epsilon_k.
\end{equation}

The fact that we use $\theta_k$ instead of $\theta$ is an important difference. The random-effects model only assumes that $\theta_k$ is the true effect size of one single study $k$. It stipulates that there is a second source of error, denoted by $\zeta_k$. This second source of error is introduced by the fact that even the true effect size $\theta_k$ of study $k$ is only part of an over-arching distribution of true effect sizes with mean $\mu$ \autocite{Harrer2021}.

\begin{equation}
\theta_k = \mu + \zeta_k.
\end{equation}

The random-effects model tells us that there is a hierarchy of two processes. The observed effect sizes of a study deviate from their true value because of the sampling error. But even the true effect sizes are only a draw from a universe of true effects, whose mean $\mu$ we want to estimate as the pooled effect of our meta-analysis. We can express the random-effects mode in one line. This formula makes it clear that our observed effect size deviates from the pooled effect $\mu$ because of two error terms, $\zeta_k$ and $\epsilon_k$ \autocite{Harrer2021}.

\begin{equation}
\hat{\theta}_k = \mu + \zeta_k + \epsilon_k.
\label{Equation:RandomEffects}
\end{equation}

A crucial assumption of the random-effects model is that the size of $\zeta_k$ is independent of $k$. That is, the size of $\zeta_k$ is a product of chance, and chance alone. This is known as the exchangeability assumption of the random-effects model. All true effect sizes are assumed to be exchangeable in so far as we have nothing that could tell us how big $\zeta_k$ will be in some study $k$ before seeing the data \autocite{Harrer2021}.

\begin{figure}[!htbp]
\centering
\includegraphics[width=0.8\linewidth]{appendix/methods/analysis/assets/ma/random_effects.png}
\caption[Random-effects model]{Illustration of parameters from the random-effects model.}
\end{figure}

The challenge associated with the random-effects model is that we have to take the error $\zeta_k$ into account. To do this, we have to estimate the variance of the distribution of true effect sizes. This variance is known as $\tau^2$. Once we know the value of $\tau^2$, we can include the between-study heterogeneity when determining the inverse-variance weight of each effect size. In the random-effects model, we therefore calculate an adjusted random-effects weight $w_k^\ast$ for each observation \autocite{Harrer2021}.

\begin{equation}
w_k^\ast = \frac{1}{s_k^2+\tau^2}.
\end{equation}

Using the adjusted random-effects weights, we then calculate the pooled effect size using the inverse variance method, just as we do using the fixed-effect model \autocite{Harrer2021}.

\begin{equation}
\hat{\theta} = \frac{\sum_{k=1}^{K}\hat{\theta}_kw_k^\ast}{\sum_{k=1}^{K}w_K^\ast}.
\end{equation}

There are several methods to estimate $\tau^2$, and it is an ongoing research question which of these estimators performs best for different kinds of data. Overall, estimators of $\tau^2$ fall into two categories. Some, like the DerSimonian-Laird and Sidik-Jonkman estimator, are based on closed-form expressions. The restricted maximum likelihood, Paule-Mandel estimator, and empirical Bayes estimator find the optimal value of $\tau^2$ through an iterative algorithm \autocite{Harrer2021}.

The Knapp-Hartung adjustments try to control for the uncertainty in our estimate of the between-study heterogeneity. While significance tests of the pooled effect usually assume a normal distribution known as Wald tests, the Knapp-Hartung method is based a $t$-distribution. Applying a Knapp-Hartung adjustment is usually sensible. Several studies have shown that these adjustments can reduce the chance of false positives, especially when the number of studies is small \autocite{Harrer2021}.
\end{document}
